\documentclass[11pt]{article}
\usepackage[a4paper,margin=24mm]{geometry}
\usepackage[T1]{fontenc}\usepackage{lmodern}
\usepackage{amsmath,amssymb,amsthm,booktabs,longtable,array}
\usepackage{microtype,xcolor,hyperref,fancyvrb,fvextra}
\hypersetup{colorlinks=true,urlcolor=blue!55!black,linkcolor=blue!45!black}
\DefineVerbatimEnvironment{Highlighting}{Verbatim}{breaklines,fontsize=\small}
\newcommand{\Con}{\operatorname{Con}}
\newcommand{\Z}{\mathbb Z}\newcommand{\N}{\mathbb N}
\newtheorem{theorem}{Theorem}\newtheorem{lemma}{Lemma}
\setlength{\emergencystretch}{3em}
\title{Joint positivity, free Lucas coordinates, and commutative proof certificates\\\large A tenth synthesis of explicit Diophantine consistency constructions}
\author{Research packet for the ProveIt repository}
\date{23 October 2026}
\begin{document}
\maketitle
\begin{abstract}
This iteration separates two kinds of progress. First, a nearest-square
identity compresses two strict-positivity requirements over unrestricted
integers into one. Applied to the pinned source, it gives a single
polynomial in 28 integers of degree eight, and a companion in 26 integers
of degree sixteen. A second completed reduction removes a high-degree
Lucas-coordinate parameterization by proving that its usual congruence
is redundant on the actual source slice. Second, an alternative proof
source dispenses with the ordered stack: single-conclusion derivability
is characterized by a finite nonnegative flow, then by products of
pairwise coprime affine integers. Its local polynomial is explicit,
but its global Diophantine compilation is not yet materialized or counted.
An asymmetric three-Pell parity interface is also proved independently.
Every final claim distinguishes a completed consistency equation from
a local module or a proposed compiler. No global minimality, new
Lean/Rocq certificate, or historical priority is claimed.
\end{abstract}
\setcounter{tocdepth}{1}
\tableofcontents
\clearpage
\section{Scope, source pin, and recommended reading}
The target is an ordinary polynomial with fixed integer coefficients
whose absence of solutions is provably equivalent, in ZFC, to
$\Con(\mathrm{ZFC})$. In particular, no coefficient may depend on whether
ZFC is consistent. Huge fixed integers are acceptable, but are part of
the specification rather than hidden parameters. The artifacts include
actual coefficient payloads, exact arithmetic circuits, mathematical
proofs, source bindings, and independently authored finite checks.

The reviewed repository is
\url{https://github.com/VladimirReshetnikov/ProveIt}, main commit
\texttt{9cca523d9f27ad39d52993f4065cbd7e58f8bdd4}.
Eight newly integrated packages, source60 through source67, were reviewed
as inert data. They were placed by commit
\texttt{a2cd79b9120fc75f48dcbef6e38d0f8b18d64d49}.
The detailed archive/member hashes and source paths appear in
\texttt{review/SOURCE\_PINS.json} and \texttt{PROVENANCE.json}.
No program from an incoming archive was executed.

\paragraph{What ``simplest'' means here.}
There is no single scalar objective. We keep the unknown domain, number
of unknowns, total degree, and binary arithmetic-operation count separate.
A square is a multiplication in the circuit ledger; named registers are
abbreviations, not additional quantified variables. An unrestricted
integer count is not interchangeable with a positive or nonnegative
count. Local witnesses in a bounded computation are not the variable
count of the fully compiled equation.

\paragraph{Logical status.}
A completed equation below has a mathematical equivalence proof, using
explicitly named inherited source-calibration and Pell theorems. The
new finite audits establish algebraic identities, literal counts, payload
integrity, and regressions. They do not discharge the universal theorems
inside a proof-assistant kernel. Nor does an equivalence to consistency
prove consistency itself.

\section{Comparison of the eight incoming constructions}
\begin{center}\small
\begin{tabular}{@{}p{29mm}p{45mm}p{71mm}@{}}\toprule
Package & Representative completed frontier & Main reusable idea\\\midrule
60: Positive Gap & $9$ positive, degree $1{,}797{,}796$, 465 operations; $15$ positive, degree24 & Positive gap coordinates; cofinal padding; sharper carry bounds.\\
61: Congruence Carriers & $13$ positive, degree34, 93 operations; $37$ integers, degree6; $34$ integers, degree8 & Move congruences into already existing carriers rather than add witnesses.\\
62: Sign and Norm & $42$ integers, degree4; $35$ integers, degree6; $30$ integers, degree8 & Square-class separation and sign recovery; positivity guards suitable for joint compression.\\
63: Sign Recovery & $42$ integers, degree4; $14$ positive, degree34, 87 operations & Remove explicit signs by arithmetic information in existing factors.\\
64: Structural Reductions & $41$ integers, degree4; $36$ integers, degree6; $31$ integers, degree8 & Delayed square substitution and a product positivity guard.\\
65: Shared Taps & $9$ positive, degree $1{,}675{,}396$, 458 operations; $10$ integers, degree $7{,}428{,}120$, 525 operations & Equal complete coefficient records share kernel positions; height $1873$.\\
66: Previous iteration & $9$ positive, degree $1{,}845{,}940$, 444 operations; $10$ integers, degree $3{,}507{,}876$, 494 operations & Gap congruences and shorter circuits; one-stack local proof source.\\
67: Oriented Guards & $9$ positive, degree $1{,}494{,}006$, 487 operations; $10$ integers, degree $3{,}327{,}556$, 519 operations & Positivity-first orientation, signed square separation, suitable ten-integer finalizer.\\\bottomrule
\end{tabular}\end{center}
These are selected tradeoffs, not a claim that each dominates every other
version in its package. In particular the incoming best quartic remains
41 unrestricted integers and the incoming best sextic remains 35.
Simply combining source64's elimination with source62's elimination
cannot save another variable: the needed square row is absent in the
41-coordinate source. Counting that nonexistent row would be invalid.

The genuinely new complete reduction is the joint positivity lemma:
it changes how inequalities are represented, not the coefficient placement
or addition chain. The free-coordinate reduction is likewise structural:
it removes a source-dependent divisibility obligation rather than merely
encoding it more cheaply. The commutative proof architecture is a more
radical change of source, but its final global cost is deliberately left
unreported until compilation is completed.

\paragraph{New completed low-degree tradeoffs.}
The joint-positivity reduction takes the incoming 30-integer octic to
28 integers at the same degree. Nesting it gives 26 integers at degree16.
These preserve all arithmetic source residuals and all four fixed ports.
They do not lower the existing quartic or sextic records. The free-Lucas
construction keeps ten unrestricted integers; its primary gain is degree,
not a claim to beat every incoming operation-count record. The archived
literal circuits provide the exact source-specific operation accounting.

\begin{center}
\begin{tabular}{@{}lrrp{43mm}@{}}\toprule
New completed equation & Integer unknowns & Exact degree & Live operations\\\midrule
Joint positivity & 28 & 8 & residual-system specification\\
Nested joint positivity & 26 & 16 & residual-system specification\\
Free Lucas, shared-gap source & 10 & 1,576,864 & 481\\\bottomrule
\end{tabular}
\end{center}
The first two are explicitly given sums of squared residuals; no operation
count is assigned to an arbitrary expansion or evaluation order.

\section{A nearest-square guard: 28 unrestricted integers at degree eight}
The two independent positivity guards can be compressed without changing any
Pell equation, compiler coefficient, or arithmetic register.
For integer $x,y$, put
\[
 G(x,y,t)=xy-(t^2-x-y+1)^2.
\]
Then $x,y>0$ if and only if $G(x,y,t)>0$ for some integer $t$.
Indeed, $G>0$ forces $xy>0$, while $x=-a,y=-b$ with $a,b>0$
would give $(t^2+a+b+1)^2>ab$. Conversely, let $N=x+y-1$.
Then $xy\geq N\geq1$. Writing $m=\lfloor\sqrt N\rfloor$, choose
$t=m$ if $N-m^2\leq m$ and $t=m+1$ otherwise. The square error
has absolute value at most $m$, and its square is strictly less than $N$
(the exact-square case has error zero). Thus $G>0$.

Every positive integer is
$\operatorname{Pos}(p,q,r)=1+p^2+q^2+r(r+1)$ by the three-square theorem
applied to $4n-3$. Therefore the single residual
\[
 xy-(t^2-x-y+1)^2-1-p^2-q^2-r(r+1)
\]
replaces two separate positive ternary forms, using four private integer
coordinates rather than six. In the reviewed 30-integer octic, take
$x=F\alpha$ and $y=F(\beta J+1-F-\alpha)$.
Both are quadratic, so the new residual is quartic. All other residuals
remain unchanged, including the guard on $\eta\zeta$.
The resulting sum of 14 squared residuals has 28 unrestricted integer
variables and degree exactly eight; the unchanged coefficient
$[a^4f^4]=1$ certifies exactness.

The equivalence is existential, not a bijection of witness tuples. From
an old zero choose the nearest-square coordinate and a positive ternary
representation of $G$. From a new zero restore independent positive
ternary representations of $x$ and $y$. Composing these maps with the
pinned source theorem gives the ZFC-provable equivalence between absence
of integer zeros and $\operatorname{Con}(\mathrm{ZFC})$.
No global minimality, external priority, or new proof-assistant certificate
is asserted.

A second application gives a higher-degree companion. With
$z=\eta(1-V+jc+hE-\eta)$ and $g=G(x,y,t)$, replace all three old
positive ternary guards by
\[
 gz-(u^2-g-z+1)^2=\operatorname{Pos}(p,q,r).
\]
Applying the proved lemma twice shows that this is equivalent to
$x,y,z>0$. Nine private coordinates become five, leaving 26 integers
and 13 residuals. The new residual has coefficient $-1$ at $t^8$,
so its square has coefficient $1$ at $t^{16}$; the other residuals do
not involve $t$. The exact total degree is therefore sixteen.

\paragraph{A general joint-positivity construction.}
For any fixed $n\ge1$, put $P_1=x_1$ and recursively
\[
 P_{k+1}=P_kx_{k+1}-(t_k^2-P_k-x_{k+1}+1)^2.
\]
Repeated application of the two-input lemma gives
\[
 (x_1>0\wedge\cdots\wedge x_n>0)
 \iff\exists t_1,\ldots,t_{n-1}\in\Z\quad P_n>0.
\]
A final equation $P_n=\mathrm{Pos}(p,q,s)$ therefore uses only $n+2$
private integer witnesses instead of $3n$. The price is increasing
degree, potentially doubling at each nesting. The 28- and 26-integer
constructions are the two- and three-input instances. This is a reusable
consequence of the proved lemma, not a global optimality statement.
The total $2n+2$ integer coordinates for $n$ original positive coordinates
is consistent with the classical conversion attributed to Sun (1992) in
Sun's 2026 introduction. The claim here is the explicit nearest-square
application to these existing private guards, especially preservation of
degree eight in the 28-variable construction, not discovery of that general
variable-count conversion.

\section{A fully specified low-degree equation}
For convenient paper use, the following is a compact presentation of the
28-integer polynomial, equivalent term by term to the literal JSON. All
letters in the list
\[
 J,F,\alpha,z,f,h,r,v,\tau,\eta,y,\gamma,o,j,E,c,a,V,I,u,\rho,
 t,p,q,s,p_3,q_3,s_3
\]
are unrestricted integer unknowns. There are 28. Define abbreviations
\[
 Q=J\beta,\quad R=8\rho^2+(3+4a)\gamma,\quad
 x=F\alpha,\quad x'=F(1+Q-F-\alpha),
\]
\[
 x''=\eta(1-V+jc+hE-\eta),\qquad
 \mathrm{Pos}(p,q,s)=1+p^2+q^2+s(s+1).
\]
The following twelve residuals are retained from the source:
\begin{align*}
 R_1={}&8\rho^2-4u^2(1+Q),\\
 R_2={}&a-E-(1+2v)(1+Q)^3,\\
 R_3={}&E+8\rho^2(E-a),\\
 R_4={}&c-\eta+(a-E)(V-1-jc-hE),\\
 R_5={}&\tau^2-1-E(c-\eta)(1-V+jc+hE)-E^2(c-\eta)^2,\\
 R_6={}&R^2+2caR-(3+4a)c^2-1,\\
 R_7={}&-V+jc-JM_C-JM_F(1+Q)\\
       &+(Q+2)\bigl(FQ^2-\alpha Q-Q^2(Q+1)\bigr),\\
 R_8={}&(1-\alpha-F+Q)(K_0+4u^2)-F-Qz,\\
 R_9={}&V+c-fo,\\
 R_{10}={}&I-(1+2r)c^2,\\
 R_{11}={}&I^2+(1-f^2)(a^2+4a+3),\\
 R_{12}={}&-1+V^2I^2+y^2-y^2I^2.
\end{align*}
Add two residuals
\[
 R_{13}=xx'-(t^2-x-x'+1)^2-\mathrm{Pos}(p,q,s),\qquad
 R_{14}=x''-\mathrm{Pos}(p_3,q_3,s_3).
\]
The promised single ordinary equation is
\[
 \boxed{\quad P_{28}=\sum_{i=1}^{14}R_i^2=0.\quad}
\]
The letter $v$ here is the JSON variable \texttt{v\_scale}; the four
new nearest-square/positivity variables are
\texttt{pair\_root}, \texttt{pair\_pos1}, \texttt{pair\_pos2},
\texttt{pair\_pos3}. The last three are \texttt{EZ\_pos1},
\texttt{EZ\_pos2}, \texttt{EZ\_pos3}. Thus there is no implicit extra
coordinate in this compact notation.

\paragraph{Fixed coefficients.}
The four constants are the whole authenticated source34 tuple:
\[
 \beta=2^{1064921}-1,\qquad K_0=K,\quad M_C=MC,\quad M_F=MF_{\rm source}.
\]
They are fixed integers, supplied as minimal unsigned little-endian
payloads under \texttt{integer\_system\_audit/coefficients/}.
The last port is already shifted; replacing it by the native unshifted
mask is invalid. The binding gives hashes, bit lengths, the 943 ordered
tiles, and the explicit bounded coefficient recipe. The independently
authored checker reconstructs all four integers from that recipe and
compares them with the materialized bytes. The total raw payload is
532,479 bytes, or 10,466 bytes in the original compressed copies.

\paragraph{The 26-integer companion.}
Keep the first 21 variables and $R_1,\ldots,R_{12}$. Introduce only
$t,t',p,q,s$ and set $g=xx'-(t^2-x-x'+1)^2$. The equation is
\[
 \boxed{\quad P_{26}=\sum_{i=1}^{12}R_i^2+
 \bigl(gx''-(t'^2-g-x''+1)^2-\mathrm{Pos}(p,q,s)\bigr)^2=0.\quad}
\]
Its exact degree is sixteen. In ZFC one can prove, for each of these
fixed polynomials,
\[
 \neg\exists\mathbf x\in\Z^n\ P_n(\mathbf x)=0
 \quad\longleftrightarrow\quad\Con(\mathrm{ZFC}).
\]
The new proof is the explicit existential reduction to source30, not a
claim that all of source30's inherited calibration has been re-proved
in this packet. That dependence is made inspectable by retaining the
source30/Q42X residual map, the sign-recovery proofs, compiler proof,
and source34 coefficient evidence. Transitive proof-assistant imports
remain references to the pinned repository.

\newcommand{\FreeCMainDegree}{1\,576\,864}
\newcommand{\FreeCMainOps}{481}
\newcommand{\FreeCParameterDegree}{16\,424}
% Main free-coordinate construction; source-specific metrics may be redefined by assembler.
\providecommand{\FreeCMainDegree}{1\,629\,856}
\providecommand{\FreeCMainOps}{517}
\providecommand{\FreeCParameterDegree}{16\,976}
\section{A free Lucas coordinate: removing the high-degree index carrier}
\label{sec:free-lucas-coordinate}

The central new construction has ten unrestricted integer unknowns.
Its literal source-specific circuit has exact degree $\FreeCMainDegree$ and
$\FreeCMainOps$ live binary arithmetic operations. The mathematical change
is not a better addition chain or a different source layout. We prove that
a high-degree defining congruence in the Lucas certificate is unnecessary
on a fixed residue slice. Its former quotient coordinate becomes a direct,
low-degree Lucas coordinate. The degree formula changes from
$196a+260$ to $96a+160$, where $a=\deg A$.

All statements here concern ordinary mathematics formalizable in ZFC.
The native compiler/refutation equivalence, Pell classification and the
identified Lucas index lemmas remain explicit dependencies. This is not a
new Lean or Rocq certificate, a consistency decision, or a global minimality
claim.

\subsection{The literal source prefix}
The selected positive-gap shared-tap certificate has $W=116$, $H=1873$,
31 weights $w_i$, 22 targets $q_j$, and three kernel polynomials $D_j$.
The complete finite lists are in \texttt{norm/final10/data/}; their
coefficient-extraction identities are proved in the source section below.
Put $\kappa=2134940$, $b=6t+4$, $g=bg_0$, and define
\begin{gather*}
 B=2^\kappa b^2,\quad \ell=\sum_{i=1}^{31}B^{w_i},\quad
 U_c=\ell+g,\quad G_c=D_2(B)U_c^2+D_1(B)U_c+D_0(B),\\
 A_T=\sum_{j=1}^{22}B^{q_j},\quad N_0=B^{117},\quad N=B^{1991},\\
 \mathrm{mask}=N_0-1-(b-1)\ell,\quad J_c=2G_c+A_T,\\
 S_c=g+N_0J_c,\quad T_c=\mathrm{mask}+N_0(B-2)A_T-2,\\
 X=N(S_c+T_c+1)-S_c-1,\\
 Y=2^{16\kappa-23}bB^{2114},\qquad\mathcal C=bB^{116}.
\end{gather*}
Every displayed sum has a fixed finite number of terms, and every exponent
is a fixed integer. The whole native four-port tuple and the twelve
materialized circuit leaves are bound independently; they must not be
substituted by the different source34 tuple used in $P_{28}$ and $P_{26}$.
The new degree bounds are $\deg X=7962$, $\deg Y=4229$ and
$\deg A=16424$. The independently checked primary circuit uses 301
multiplications, 155 additions and 25 subtractions, totaling 481.
Compared with the incoming ten-integer degree-$3{,}327{,}556$/519-operation
construction, both metrics decrease. It also improves the previous
iteration's 494-operation, degree-$3{,}507{,}876$ variant.
The new literal circuit treats the materialized threshold integer
$2^{34159017}$ as one fixed leaf. Retaining the predecessor's five
constant-building gates would give 486 operations, not 481. Both counts
are reported to separate structural simplification from constant folding;
neither is a bit-complexity or execution-time estimate.

\subsection{The complete finalizer}

Take the fixed source polynomials $X,Y,\mathcal C$ and its normalizations
\[
 f_{\rm old}=2t+1,\qquad b=6t+4,\qquad
 g=bg_0,\qquad k=b(3k_0+2).
\]
The exact source prefix and all four coefficient ports are specified in the
accompanying source section and literal data. Its universally valid
congruences are
\[
 X\equiv1\pmod2,\quad X\equiv0\pmod3,\quad
 b\mid X+1,\quad b\mid Y.
\]
The ten supplied variables are
$(t,g_0,c_0,k_0,l,w,x,y,z,m)\in\mathbb Z^{10}$. Every other symbol below
is a polynomial abbreviation. Define
\begin{gather*}
 L=lY,\quad U=2LX,\quad V=4gwY,\quad W=bw,\\
 A_{1/2}=(LX)(V+1),\quad A=2A_{1/2},\quad B_1=2X+1,
 \qquad \boxed{C=6c_0-1},\\
 P_{1/2}=(LX)UV,\quad P=2P_{1/2},\quad
 K=X+1+k(P-2),\\
 S=2A-5,\quad T=3WC-2(W^2-1),\quad
 Q=(P_{1/2}^2-1)K^2+1,\\
 \Delta=8\mathcal C^3gK^2-g^2
       \bigl(32(C-KL)^2\mathcal C^3+g^2K^2\bigr),\\
 \xi=zx,\quad D=(A^2-4)C^2+4,\quad E=C^2D\xi,\\
 F=4(A^2-4)E^2+1,\quad
 G=2(A+2)(A-2)^2E^2-1,\\
 \Omega=16E,\quad H=C+F(B_1-C+\Omega y),\quad
 I=(G^2-1)H^2+1.
\end{gather*}
The circuit uses the identical expression
$G=(A_{1/2}-1)(F-1)-1$ to share operations. The subscripts $1/2$ name
explicit polynomial registers, not divisions in the equation.

Keep the original one-modulus intersective finalizer:
\begin{gather*}
 \rho=T+Sy,\qquad \Sigma=\xi S^5,\\
 \Theta=(\rho^3-\rho S^2-S^3)(\rho^2+23S^2),\\
 r=bQ,\qquad \delta=9\Delta-3,\qquad
 a_0=r+\delta(rz)^2,\quad a_1=DFI,\quad Z=\Sigma m-\Theta.
\end{gather*}
The single equation is
\begin{equation}\label{eq:free-C-final}
 \boxed{\bigl(Z^2-\Sigma^2(a_0+a_1)\bigr)^2
             -4\Sigma^4a_0a_1=0.}
\end{equation}
Every exponent is fixed. There are no hidden quotient, root, binomial,
variable-exponent, or bounded-universal operations. All fixed leaves are
materialized integers with hashes. In particular the new parameterization
has no requirement $(A-2)\mid C-B_1$ built into it.

\subsection{A strengthened signed Lucas certificate}

\paragraph{Lemma.}
Suppose $A,B,C,\xi,y\in\mathbb Z$ satisfy
\[
 6\mid A,\quad B\equiv3\pmod4,\quad
 1<B<|A|/2-1,\quad C\equiv5\pmod6,\quad\xi\ne0.
\]
Use $D,E,F,G,\Omega,H,I$ as above with $B$ in place of $B_1$.
If $DFI$ is a square, then $C=u_B(A)$, where
$u_0=0,u_1=1,u_{n+2}=Au_{n+1}-u_n$.
No hypothesis $(A-2)\mid C-B$ is needed.

\paragraph{Proof.}
Here $C$ is odd and nonzero. Since $|A|>2$, the three integers $D,F,I$
are positive. The certificate identities
\[
 F\equiv I\equiv1\pmod D,\qquad 4I\equiv D\pmod F
\]
and oddness of $F$ make them pairwise coprime, so they are separately
squares. The usual Pell classification gives integers $s,\mu,t$ with
\[
 C=u_s(A),\qquad 4E=u_\mu(A),\qquad H=u_t(2G).
\]
Put $n=|\mu|>3$. The parity identities show that $n$ is even and $s,t$
are odd. The Lucas index-separation lemma applied to the congruence modulo
$F$ gives $s\equiv\pm t\pmod{2n}$. The rank-divisibility lemma applied
to $C^2\mid u_n(A)$ gives $C\mid n$. Finally $H\equiv B\pmod\Omega$
and $2G\equiv-2\pmod\Omega$ give $t\equiv B\pmod\Omega$.
Thus
\begin{equation}\label{eq:free-C-rank-residue}
 s\equiv B\pmod{2C}\quad\hbox{or}\quad s\equiv-B\pmod{2C}.
\end{equation}
These are precisely the congruence-free classification and index steps of
the modified certificate; the new work is to remove its exceptional ranks
and sign without the old defining congruence.

For every $j\ge1$, direct paired induction in the recurrence proves
\begin{equation}\label{eq:free-C-paired-recurrence}
 u_{2j}(A)\equiv(-1)^{j-1}jA\pmod{A^2},\qquad
 u_{2j+1}(A)\equiv(-1)^j\pmod{A^2}.
\end{equation}
The base is $u_2=A,u_3=A^2-1$. The next two recurrence steps multiply the
second congruence by $A$ and subtract the first, then multiply again by $A$
and subtract the second. This works for negative $A$ as well.

If $s=\pm1$, then $C=\pm1$ and $D=A^2$. Therefore $A^2\mid u_n(A)$.
Writing $n=2j$, the first congruence in
\eqref{eq:free-C-paired-recurrence} yields $A\mid j$.
Consequently $2|A|$ divides both $2n$ and
$\Omega=16A^2\xi$. Combining $s\equiv\pm t\pmod{2n}$ with
$t\equiv B\pmod\Omega$ forces $B\equiv\pm1\pmod{2|A|}$,
contradicting $1<B<|A|/2-1$.
This argument does not assume the generally false divisibility $n\mid u_n(A)$.

Hence $|s|\ge2$. Lucas growth gives
$|C|=u_{|s|}(|A|)\ge|A|>B$ and $|s|\le|C|$.
The two congruences in \eqref{eq:free-C-rank-residue} therefore force
$s=B$ or $s=-B$: a nonzero multiple of $2C$ would be too large.
Since $B\equiv3\pmod4$, the second congruence in
\eqref{eq:free-C-paired-recurrence} gives $u_B(A)\equiv-1\pmod{A^2}$,
and thus modulo3. The branch $s=-B$ would give $C\equiv1\pmod3$,
contrary to $C\equiv5\pmod6$. Therefore $s=B$, as required.
\hfill$\square$

The odd, nonzero residue slice is essential: $A=12$, $B=3$, $C=0$, $\xi=1$, $y=0$ gives $DFI=4$, although $u_3(12)=143$.

\paragraph{A reusable unsigned interface.}
The mechanism is not restricted to this particular circuit. For even $A$,
odd $C$, odd $B$ with $1<B<|A|/2-1$, and nonzero $\xi$, the same
square-certificate index argument gives $C=\pm u_B(A)$ without assuming
$A-2\mid C-B$. To handle both exceptional indices $s=\pm1$, use
$C^2=1$: then $D=A^2$, $4E=4A^2\xi$ and the paired recurrence forces
$2|A|\mid n$, just as above. Both exceptions contradict the range of $B$.
For $|s|\ge3$, the growth and index congruences give $s=\pm B$.
The fixed residue conditions $3\mid A$, $B\equiv3\pmod4$,
$C\equiv-1\pmod3$ then select the positive index. Thus the primary
construction is a corollary of a stronger arithmetic interface, not merely
a shorter evaluation DAG. Its dependency on the standard Lucas index
comparison is unchanged; no claim of literature priority is made.

\subsection{Signed soundness of the ordinary equation}

The source and $k$ normalization imply $b\mid K$, $3\mid K$, and
\[
 Q\equiv1\pmod{b^2},\qquad Q\equiv1\pmod3,\qquad \gcd(b,Q)=1.
\]
On every signed assignment $3\mid A$, $C\equiv-1\pmod3$, so $D\equiv0$
modulo3 and $a_1\equiv0\pmod3$. In contrast $r=bQ\equiv1\pmod3$ and
$\delta\equiv0\pmod3$, so $a_0\equiv1\pmod3$. Thus
\emph{$a_0\ne a_1$ unconditionally}, before any positivity argument.

The odd integer $S=2A-5$ is nonzero. The polynomial
$\pi(v)=(v^3-v-1)(v^2+23)$ has no rational root; hence
$\Theta=S^5\pi(\rho/S)\ne0$. A zero of
\eqref{eq:free-C-final} therefore has $\Sigma\ne0$. With $q=Z/\Sigma$,
the equation gives $(q^2-a_0-a_1)^2=4a_0a_1$.
Since $a_0\ne a_1$, $q\ne0$. The rationals
\[
 u=\frac{q^2+a_0-a_1}{2q},\qquad v=q-u
\]
have squares $a_0,a_1$. Rational-square-root integrality makes $u,v$
integers. Thus $a_0,a_1$ are squares and $\Sigma\mid\Theta$.
A prime-valuation comparison in the homogeneous quintic yields $S\mid T$:
if $v_p(T)<v_p(S)$, the two factors of $\Theta$ have valuations
$3v_p(T)$ and $2v_p(T)$, contradicting $S^5\mid\Theta$.

The factorization $a_0=r(1+\delta r z^2)$ has two coprime nonzero factors.
Since $r\equiv1\pmod3$, the signed-square possibility makes $r$ and the
bracket positive squares. Applying the same argument to $r=bQ$ gives
$b,Q$ positive squares, hence $b\ge4$. If $\Delta\le0$, then
$\delta=9\Delta-3\le-3$, making the bracket negative. Thus
$\Delta>0$. Its defining inequality gives $K\ne0$ and $0<g<2\mathcal C$.
If $l=0$ then $Q=1-K^2\le0$, a contradiction. Hence $l\ne0$;
$f_{\rm old}=2t+1$ is already nonzero.

The source bounds now give $X>3b$ and $Y>256$. Since $V+1$ is odd and
nonzero,
\[
 |A|=2|l|XY|V+1|\ge2XY>2B_1+2.
\]
All hypotheses of the strengthened lemma hold: $6\mid A$,
$B_1\equiv3\pmod4$, $C\equiv5\pmod6$, and $\xi\ne0$.
It gives $C=u_{B_1}(A)$. The recurrence modulo $A-2$ then gives
$(A-2)\mid C-B_1$ \emph{as a consequence}. Restore the proof-only integer
$h=(C-B_1)/(A-2)$ and apply the unchanged reverse binomial bridge.
The calibrated source decodes a positive source zero and hence a finite
ZFC refutation. No sign of an arbitrary supplied auxiliary variable was
silently assumed.

\subsection{Completeness}

Start from a positive source zero and its canonical arithmetic witnesses.
They have $C=u_{B_1}(A)$, $6\mid A$, and $B_1\equiv3\pmod4$.
Thus $C$ is odd and $C\equiv-1\pmod3$, giving the integer
$c_0=(C+1)/6$. The old constructive congruences still supply $k_0$.
Here $b,Q$ are positive squares and $\Delta>0$.
The positive integer
\[
 \delta=3(3\Delta-1)
\]
has $3$-adic valuation exactly1, and is therefore nonsquare; the same holds
for $\delta r$ since $r=bQ$ is a positive square.
Pell's theorem supplies positive $u,z$ with
$u^2-\delta r z^2=1$, making $a_0=ru^2$ square.

The modified Lucas interpolation theorem first supplies a positive $\xi$
divisible by this $z$, and then permits every residue class of $y$ modulo
that fixed $\xi$ while retaining $DFI$ square. Put $x=\xi/z$.
Since $S\mid T$, write $v_0=T/S$. The quintic $\pi$ has a root modulo every
positive integer, so prescribe $y$ with $\pi(v_0+y)\equiv0\pmod\xi$.
Then $\Sigma\mid\Theta$. Taking
\[
 m=\Theta/\Sigma+\sqrt{a_0}+\sqrt{a_1}
\]
gives an integer solution of \eqref{eq:free-C-final}.
Thus the construction is existence-equivalent to the fixed source, and
negating the two existence statements gives the required equivalence with
$\operatorname{Con}(\mathrm{ZFC})$.

\subsection{Exact degree and verification boundary}

Let $a=\deg A$. Since $C$ now has degree1,
\begin{align*}
 \deg D&=2a+2,& \deg E&=2a+6,& \deg F&=6a+12,\\
 \deg G&=7a+12,&\deg H&=8a+19,&\deg I&=30a+62,\\
 \deg a_1&=38a+76,&\deg\Sigma&=5a+2.
\end{align*}
For the included compatible source, the guard and center have strictly
lower degree than the dominant $\Sigma^4a_1^2$ term. Consequently
\[
 \deg F_{10}=4(5a+2)+2(38a+76)=96a+160.
\]
Exact nonzero diagonal leading coefficients are checked with the actual
fixed tuple modulo three primes. They certify that the upper degree is
attained; they are not sampled estimates of an unproved upper bound.
Separate checks authenticate coefficient bytes, exact source extraction,
full modular formulas, signed residue separation and live circuit topology.
Finite checks do not replace the unbounded mathematical proof or its
stated arithmetic/compiler dependencies.

\section{A compatible positive-gap shared-tap source}
\label{sec:gapshared}
This auxiliary improvement combines the source65 shared coefficient taps with
the round9 positive-gap coordinate substitution.  It does not change the four
native coefficients or claim a new universal-polynomial theorem.
All31 source variables range over $\mathbb Z_{>0}$.
The substitutions $P_{\rm old}=P+\eta+\zeta$ and $N_{\rm old}=N+1$ are
bijections on the positive zero sets: the original equations imply
$P_{\rm old}-\eta-\zeta=E(c-\eta)>0$ and $N_{\rm old}=q^3\ge8$.
The original31 coefficient tuple is retained together and materialized.

Let $w_i$ be the31 distinct positive weights, $q_j$ the22 targets and
$U(T,x)=\sum_i x_iT^{w_i}$.  The serialized certificates prove the formal
identities
\[
 [T^{q_j}]\bigl(D_2(T)U(T,x)^2+D_1(T)U(T,x)+D_0(T)\bigr)=2\rho_j(x)
 \quad(1\le j\le22).
\]
Equal complete coefficient records may share one kernel position; they are
not added together.  Exhaustive symbolic extraction checks every monomial,
including the absent ones.  The primary source uses86 records and81 unique
positions, compared with source65's89 records and84 positions.
Its layer degrees are $(1384,1732,1641)$, with $W=116$, $Q=1841$ and
$H=1873$.  Thus the packing removes three positions at unchanged height.
The85-record alternative does not satisfy this layout and is not substituted
silently for the primary source.

The native kernel coefficient bound remains $M=2^{2134927}$, checked on the
actual fixed integers.  For $B\ge4096Mb^2$ and
$U=\sum_i B^{w_i}+g\ge B^{116}$, the quadratic leading coefficient is1.
Its negative tail is less than half its leading term.  The possible negative
linear and constant contributions, relative to $B^{1641}U^2$, are at most
$2M/B^{25}$ and $2M/B^{489}$.  Consequently the source convolution is
positive before digit recovery.  The margin $H-Q=32$ preserves the existing
carry bounds.  Every weight is positive, so canonical encoded witnesses
remain divisible by $B$, and in particular by $b^2$.

The artifact auditor verifies110 exact target identities across five
certificates, rejects five coefficient mutations, validates all four
compressed and raw native coefficient hashes, and establishes that every
materialized kernel coefficient has absolute value below$M$.
These checks are finite algebraic certificates, not a new kernel-checked
proof of the inherited source-to-consistency theorem.

\section{Disjunctive scratch-space reuse}
\label{sec:branchscratch}
This section gives a fully explicit improvement to the \emph{local} proof
transition relation. It does not claim that a complete consistency equation
with the following variable counts has been constructed. The remaining
bounded-universal elimination is a separate cost.

\paragraph{The elementary principle.}
Suppose finitely many relations have presentations
\[
 R_j(\mathbf x)\iff\exists\mathbf z_j\in\mathbb N^{m_j}\quad
 Q_j(\mathbf x,\mathbf z_j)=0.
\]
Put $m=\max_j m_j$. Then
\[
 \bigvee_j R_j(\mathbf x)\iff
 \exists\mathbf z\in\mathbb N^m\quad
 \prod_j Q_j(\mathbf x,z_1,\ldots,z_{m_j})=0.
\]
Indeed a product of integers is zero exactly when a factor is zero. Only
one branch's scratch values need be correct; unselected branches impose
no obligations. In contrast, flattening the whole arithmetic circuit
before taking its zero set generally forces scratch equations for every
branch and pays a sum rather than a maximum.

\paragraph{Application to the one-stack source.}
Use $I(a,b)=2(a+b)(a+b+1)+4b$, with membership, universal
quantification, and equality obtained by adding $1,2,3$ respectively.
Let $L(s,h,t)=2s-2-(h+t)(h+t+1)-2t$.
The equation $L(s,h,t)=0$ expresses a stack cell with head $h$ and tail $t$.
There are twelve alternatives: truth push, eight axiom-template pushes,
the guarded seed-bundle push, modus ponens, and generalization. Each
branch has its own arithmetic circuit, but the variable names
$q_0,q_1,\ldots$ are reused across all branches.
Each distinct constructor uses one nonnegative wire
\[
 q=(u+v)(u+v+1),\qquad I(u,v)=2q+4v.
\]
The four shared parameters are $a,b,c,d$. In the eighth axiom $d$ is the
fourth schema parameter. In modus ponens it is the intermediate stack;
in the bundle it certifies the guard. These roles need not agree between
unselected alternatives.

For the first construction the bundle guard is
\[
 q_{26}=(c-a)^2,\quad q_{27}=(c-b)^2,\quad
 q_{26}q_{27}=(d+1)^2.
\]
All residuals are quadratic. The twelve sums of their squares have degree
four, giving a product of degree $48$, with $28$ common wires and four
parameters: $32$ local witnesses, compared with the previous $89$.
The exposed stack coordinates are not included in these counts.

\paragraph{The recommended local tradeoff.}
Instead use the literal guard
\[
 ((c-a)(c-b))^2-(d+1)^2=0.
\]
Over natural parameters this holds for some $d$ precisely when $c\ne a,b$.
It saves the two square wires. Also substitute the defining polynomial of
the bundle's outermost constructor wire into its stack-cell residual and
omit that wire equation. The substitution is exact and does not alter
existential solvability. The bundle sum of squares now has degree eight;
the eleven remaining factors still have degree four. Thus the product has
exact degree $8+11\cdot4=52$ and only $25+4=29$ local witnesses.
No expansion of the twelve-factor product is necessary: the attached
JSON lists every residual and each expanded branch sum of squares, and
specifies that their product is the final local polynomial.

\paragraph{Correctness.}
If the polynomial vanishes, one branch sum of squares vanishes, and hence
each of its residuals vanishes. By induction through that branch's
constructor circuit its wires have their intended values. The stack
residuals therefore state the corresponding legal proof transition. The
bundle guard gives the required disequalities. Conversely, for any legal
transition choose its schema parameters and assign the selected branch's
wires their ordinary nonnegative values; choose arbitrary values for
unused common wires. All selected residuals vanish. Eliminating an
outermost wire preserves both directions because its value is uniquely
determined and nonnegative. This proves equality of the local relations,
not merely agreement on tested examples.

\paragraph{Limits and reproducibility.}
The JSON files are literal ordinary integer polynomials, with no hidden
variable exponents or unbounded products: the outer product has exactly
twelve factors. The global equivalence to $\neg\operatorname{Con}(\mathrm{ZFC})$
would follow by substituting this local relation in the established finite
one-stack-run characterization and completing bounded-universal
elimination. That global constructor is not executed here, so these local
counts must not be compared directly with the existing nine- and
ten-variable completed consistency equations. The independent checker
reads the JSON without importing the generator, evaluates selected
branches using a separate integer interpretation of the constructors,
and rejects endpoint mutations. These tests supplement the algebraic
argument; they are not kernel certificates.

\paragraph{Dependency-aware refinement: seventeen local witnesses.}
The previous paragraph can be improved without increasing the degree.
Introduce a wire for the \emph{whole} constructor output
$z=I(u,v)=2(u+v)(u+v+1)+4v$, rather than a wire only for
$(u+v)(u+v+1)$. Whole-output wires prevent linear tails from propagating
extra dependencies. The bundle now has $27$ constructor wires. In their
undirected dependency graph, the set
\[
 \{1,3,4,5,7,9,12,15,17,19,21,22,24,26\}
\]
is independent. Substitute the defining quadratic of each selected wire
and delete its defining residual. Since no selected wire depends on
another selected wire, the substituted expressions have degree two, and
each remaining residual has degree at most four. There are only $13$
remaining wires, with the same four base parameters. The twelve-factor
local transition polynomial therefore has $17$ witnesses and degree $52$.
Its expanded bundle sum of squares has $1110$ monomials.
The global correctness argument is unchanged: this is exact existential
elimination of nonnegative, uniquely defined scratch values.

For use in the proof-flow architecture, replace stack-cell residuals by
$p=q=0,\ v=\rho$ in each axiom branch, by $q=I(p,v)$ in the modus-ponens
branch, and by $q=0,\ v=\operatorname{All}(a,p)$ in the generalization
branch. The truth branch is $p=q=v=0$. The resulting literal polynomial
has three exposed formula labels, the same $17$ local witnesses, and
degree $50$: the truth factor now has degree two rather than four.
The independent-wire elimination uses the same fourteen-index set.

\section{A different source: commutative proofs and affine products}
\label{sec:commutative}
The main architectural idea in this component is to discard proof order.
The following result explains why a cycle check is unnecessary. It is a
mathematical theorem, not a claim of historical priority for flow certificates.

\paragraph{Single-conclusion flow theorem.}
Let $F$ be a set of formula codes and $R$ a family of finitary rules, each
having a premise multiset $P_r$ and exactly one conclusion $c_r$. Axioms
have empty premise multisets. Write $D$ for the least rule-closed subset of
$F$, and $e_f$ for the basis vector of $f$ in $\mathbb Z^{(F)}$. Then
\[
 t\in D\quad\Longleftrightarrow\quad
 \exists(n_r)\in\mathbb N^{(R)}\quad
 \sum_r n_r\left(e_{c_r}-\sum_{p\in P_r}e_p\right)=e_t.
\]
Completeness counts occurrences in a finite proof tree: every non-root
occurrence cancels once. For soundness, suppose $t\notin D$, and sum the
coordinates outside $D$. Each rule contributes at most zero: an
underivable conclusion must have an underivable premise. The right-hand
side contributes one, a contradiction. This argument even proves soundness
for nonnegative rational multiplicities and a coordinatewise lower bound
by $e_t$. A finite certificate can be converted constructively into a
proof DAG by saturating its finite set of selected rules.

The single-conclusion restriction is essential. The rule with premise $a$
and two conclusions $a,t$ has balance $e_t$, without deriving either
formula from no axioms. Thus the result is not a general Petri-net
reachability criterion. Nonnegative multiplicities are likewise essential.

\paragraph{Free truth and uniform arity.}
In the calibrated C222 calculus, truth code $0$ is derivable. We may
therefore ignore coordinate $0$, pad every axiom with two zero premises,
and pad generalization with one zero premise. Every rule now has the form
$(p,q)\longrightarrow v$. The target is code $1$. The inherited C222-to-ZF
calibration, followed by the constructible-choice bridge from ZF to ZFC,
shows in ZFC that existence of such a finite flow certificate is equivalent
to $\neg\operatorname{Con}(\mathrm{ZFC})$.

\paragraph{Coprime affine weights.}
Choose $m\geq1$ bounding every label in a proposed certificate, and a
positive common multiple $M$ of $1,\ldots,m$. Set
\[
 f(j)=1+jM\qquad(0\leq j\leq m).
\]
Then $f(0)=1$, and the positive-label weights are pairwise coprime.
Indeed, a common divisor of $1+iM$ and $1+jM$ is coprime to $M$ and divides
$j-i$; the latter divides $M$, forcing that common divisor to be one.
Unique factorization consequently gives the exact equivalence
\[
 \prod_{i<N}(1+v_iM)
 = (1+M)\prod_{i<N}(1+p_iM)(1+q_iM)
\]
if and only if the uniform rules have balance $e_1$ modulo coordinate $0$.
No primality predicate and no exponentiation by a formula code appears.
Taking $M=m!$ proves existence; factorial is not being treated as a polynomial.
Merely taking an arbitrary large $M$ is not a substitute for the coprimality
condition.

\paragraph{A one-accumulator source.}
A product certificate can be realized by a path
\[
 S_0=1+M,\qquad S_N=1,\qquad
 S_{i+1}(1+v_iM)=S_i(1+p_iM)(1+q_iM).
\]
Multiplying these equations and cancelling the positive internal
accumulators gives the product identity. Positivity follows inductively
from $S_0>0$ and positivity of every factor. Conversely, expand the nodes
of a proof tree, starting with its sole root obligation. Let $S_i$ be the
product of the weights of the current open obligations. Expanding a node
replaces its conclusion factor by its two premise factors and gives the
specified update. Truth obligations have factor one. At the end the
accumulator is one. The proof does not require an ordered stack, decoding
of a whole stack, or ranks on proof nodes.

\paragraph{Exact local polynomial.}
Let $H(p,q,v;\boldsymbol z)$ be the branch-shared uniform-Horn polynomial
in \texttt{horn\_shared\_source.json}. It is the product of twelve sums
of squares, is nonnegative on all integer assignments, and uses seventeen
natural local witnesses. Its exact degree is fifty. Set $M=M_0+1$ and
introduce natural label slacks $d_p,d_q,d_v$. The local accumulator
polynomial is exactly
\begin{align*}
 A={}&H(p,q,v;\boldsymbol z)\\
 &+\bigl(S'(1+v(M_0+1))-S(1+p(M_0+1))(1+q(M_0+1))\bigr)^2\\
 &+(p+d_p-m)^2+(q+d_q-m)^2+(v+d_v-m)^2.
\end{align*}
Its nonnegativity means $A=0$ enforces every displayed condition, including
one genuine Horn rule. There are $17+3+3=23$ local witnesses: the Horn
witnesses, three labels, and three slacks. The four exposed coordinates
are $S,S',m,M_0$. The update has total degree five before squaring, so the
whole polynomial still has degree fifty. This is a \emph{local} count,
not the variable count of the eventual consistency equation.

\paragraph{A closed, but uncompiled, constructor.}
Define $\beta(x,y,i)=x\bmod(1+(i+1)y)$. The five outer natural coordinates
are $t,x,y,m,M_0$. The exact input predicate for the supplied constructor is
\begin{align*}
 &\beta(x,y,0)=M_0+2,\qquad \beta(x,y,t)=1,\\
 &\forall i<t\;\exists\boldsymbol w\in\mathbb N^{23}\quad
 A\bigl(\beta(x,y,i),\beta(x,y,i+1),m,M_0;\boldsymbol w\bigr)=0,\\
 &\forall j<m\;\exists k\in\mathbb N\quad M_0+1=(j+1)k.
\end{align*}
The last line supplies the common-multiple property directly. If $m=0$,
all labels are zero and the positive accumulator cannot change, so the
endpoints exclude this degenerate case. Similarly $t=0$ is excluded.
For completeness, select $m$, then $M=m!$, then a finite accumulator path,
and finally CRT witnesses $x,y$. The two bounded blocks are converted
by the existing bounded-forall constructor; all remaining coordinates are
existentially closed; the direct quartic compiler is applied before, and
without, the Jones finalizer.

\texttt{SharedAffineProductEquation.lean} specifies this entire constructor
against the pinned APIs. It has not been elaborated or kernel checked,
and its resulting ordinary polynomial has not been materialized or
counted. The common multiple $M$, beta scale $y$, and the bounded-forall
compiler's internal CRT/majorization scales are distinct witnesses.
Identifying them could create a circular impossible bound, and is not
part of this construction. The result therefore supplies a genuinely
different source and a precise compilation target, not a claim that a
smaller final consistency polynomial has already been obtained.

\paragraph{Checks and scope.}
The direct finite checks cover $340{,}878$ target/flow cases,
$166{,}650$ pairwise-gcd cases, and $227{,}252$ product/balance comparisons.
The shared local polynomial has twelve independently re-expanded
sum-of-squares factors and $480$ modular local replays. Negative controls
exhibit the dangers of multiple conclusions, signed multiplicities,
insufficient radix, and missing coprimality. These are finite regressions,
not replacements for the unbounded mathematical proofs or a Lean check.

\section{A standalone asymmetric parity interface}
\label{sec:parity}
This component bypasses the inflated central-binomial index, but does
\emph{not} yet constitute a new complete consistency equation. Its
precise theorem concerns positive input integers $S,T$:
\[
 S\mathbin{\&}T=0\quad\Longleftrightarrow\quad
 \exists\mathbf w\in\N^{12}\ P(S,T,\mathbf w)=0.
\]
The bit operation occurs only in this specification, not in the polynomial.
The three-Pell quotient technique has a precedent in Jones--Sato--Wada--Wiens
(1976), Section3.9. The specialization here uses odd-quotient proximity,
rank locking, and a built-in congruence to remove one witness.

Write $n=S+T$ and use the twelve nonnegative witnesses
$y,r,c,p,l,i,j,u_0,u_1,u_2,u_3,v$. Define
\begin{align*}
 x&=2(y+1),&Y&=2r+1,\\
 M&=128(n+1)(x+1)(Y+1),&A&=M(x+1),\\
 C&=n+1+c,&K&=T+1+p(M-1),\\
 L&=S+1+l(Mx-1),&D&=(A^2-1)C^2+1,\\
 E&=2(i+1)DC^2,&F&=(A^2-1)E^2+1,\\
 G&=A+F(F-A),&H&=C+F\bigl(n+1+(2j+1)C\bigr),\\
 I&=(G^2-1)H^2+1.&&
\end{align*}
The five residuals are
\begin{align*}
 U_0&=u_0^2-\bigl((2n+6)^3(2n+8)(y+1)^2+1\bigr),\\
 U_1&=u_1^2-\bigl((M^2-1)K^2+1\bigr),\\
 U_2&=u_2^2-\bigl(((Mx)^2-1)L^2+1\bigr),\\
 U_3&=u_3^2-DFI,\\
 U_4&=K^2L^2-64(C-YKL)^2-v-1.
\end{align*}
Set $P=\sum_{a=0}^4U_a^2$. This is a literal ordinary polynomial of
exact total degree 504, counting the two input coordinates as variables.
The independent audit corrects an earlier draft's bound of 500.
A thirteen-witness, six-equation version retains a separate divisibility
witness and has degree368. These are distinct tradeoffs.

\paragraph{Pell conventions and hypotheses.}
Let $\psi_a(k)$ be the second coordinate of
$(a+\sqrt{a^2-1})^k$. For $a>1$, every nonnegative solution of
$(a^2-1)z^2+1=\square$ is such a coordinate. Also
$\psi_a(k)\equiv k\pmod{a-1}$,
$(2a-1)^{k-1}\le\psi_a(k)\le(2a)^{k-1}$, and
\[
 (2a)^q\left(1-\frac{q}{2a}\right)
 \le\psi_a(q+1)\le(2a)^q.
\]
The growth-square theorem applied to $U_0$ gives $x>128\cdot2^n$.
The indexed-Pell certificate of JSWW Lemma3.8 gives
$C=\psi_A(n+1)$ from $C\ge n+1>1$, $DFI=\square$, and the
prescribed two congruences of $H$. Here $F\equiv1\pmod{2C}$, so
our formula for $H$ supplies $H\equiv n+1\pmod{2C}$ and
$H\equiv C\pmod F$, with $H\ge n+1+2C$.

\paragraph{Soundness and rank locking.}
The Pell equations and congruences give
\[
 K=\psi_M\bigl(T+1+a(M-1)\bigr),\quad
 L=\psi_{Mx}\bigl(S+1+b(Mx-1)\bigr),\qquad a,b\ge0.
\]
The last residual implies $|C/(KL)-Y|<1/8$ and therefore
$C/(KL)>7/8$. If $a>0$, the lower growth bound gives
$K\ge(2M-1)^{2n+1}>2(2A)^n\ge2C$, because
$M-1\ge2n+1$ and $(2M-1)^2>2A$. The same argument applies to
$L$ if $b>0$. Either case contradicts the proximity inequality.
Thus both ranks are the intended ones.

Set $\rho=(x+1)^n/x^S$, $\sigma=C/(KL)$, and
$\epsilon=n/(2M)$. The relative Pell bounds place both normalized
numerator and denominator in $[1-\epsilon,1]$, whence
$|\rho-\sigma|\le\sigma\epsilon/(1-\epsilon)<1/128$.
The binomial expansion gives $\rho=Q+f$ with integer
$Q\equiv\binom nS\pmod x$ and $0\le f<2^n/x<1/128$.
Consequently $|Y-Q|<1/8+1/128+1/128<1$. Hence $Y=Q$ is odd.
Since $x$ is even, the binomial coefficient is odd. Lucas' binary
criterion yields $S\mathbin{\&}T=0$.

\paragraph{Completeness and the congruence compression.}
For disjoint positive $S,T$, choose a growth-square witness $y$ and
put $Y=\lfloor(x+1)^n/x^S\rfloor$, an odd positive integer. Define
$M,A,C,K,L$ at their intended Pell ranks. Their congruences supply
nonnegative $c,p,l$. The indexed-Pell theorem supplies $i$ and an
initial coordinate $H_{\rm old}$ satisfying both congruences.
Its coprimality identities separate $D,F,I_{\rm old}$ into squares,
so $H_{\rm old}$ is a Pell coordinate for base $G$.
The determinant-one Pell matrix is periodic modulo $2CF$.
Increasing its index by multiples of a positive period preserves both
congruences and gives arbitrarily large $H_{\rm new}$. Choose it at
least $C+F(n+1+C)$. Writing $H_{\rm new}-C=Fq$ yields
$q\equiv n+1+C\pmod{2C}$, hence
$q=n+1+(2j+1)C$ with $j\ge0$. This proves completeness of the
compressed formula for $H$, rather than assuming an unproved bound
for a particular canonical witness.
The same relative estimates give $|\sigma-Y|<1/64<1/8$, so the
last positive integer slack exists. The square-root witnesses exist
by construction. This completes both directions of the theorem.

\paragraph{Exact degree and boundary.}
The successive degrees of $M,A,C,D,E,F,G,H,I$ are
$3,4,1,10,13,34,68,36,208$. Thus $DFI$ has degree252 and
$U_3^2$ degree504, with a nonzero leading term; no other squared
residual cancels it. The theorem assumes $S,T>0$ and all twelve
witnesses nonnegative. Substituting a packed source would still need
a separate proved recovery of the binary base. That missing step is
not hidden in this theorem, and no completed global improvement is
claimed for this component.

\section{Verification, provenance, and proof boundaries}
\paragraph{Completed mathematical claims.}
The nearest-square lemma and the paired/nested witness transformations
are proved above over all unrestricted integer assignments. The new
Lucas slice lemma treats negative parameters and the exceptional
indices explicitly. Their composition with the pinned source theorem
gives completed consistency equations. They are not merely bounded
search results. The affine-flow theorem, disjunctive scratch reuse,
and asymmetric parity theorem are also unbounded mathematical claims,
but their new global consistency constructors are not finished.

\paragraph{Finite checks are evidence of different kinds.}
The archive checker first verifies its sealed SHA-256 file manifest and
retained-member provenance. In a fresh temporary copy it executes only
the newly authored builders and checkers. This prevents reproduction
from modifying the delivered package. Exact generated outputs must
rebuild byte for byte. Restricted AST readers are used for inherited
polynomial strings; incoming Python is never imported. The independent
source34 checker reconstructs the actual four compiler coefficients.
The direct audit re-expands every local sum of squares and checks
triangular, nonnegative constructor restoration. The displayed compact
low-degree residuals have their own exact symbolic equality check.
Finite signed, Pell-congruence, flow and modular-circuit examples are
regressions, not substitutes for the unbounded arguments.

\paragraph{Negative controls.}
The release includes a manifest tamper test and mathematical mutation
controls in the independent component audits. These have different
meanings: rejecting changed bytes establishes packaging integrity,
whereas rejecting a changed coefficient or invalid flow assumption
exercises a mathematical check. Neither should be described as a
formal kernel certificate. Optimized-Python runs are included for
independent auditors using explicit exceptions, to make clear that
those validations do not disappear under \texttt{python -O}.

\paragraph{Formal status reviewed at the pin.}
The latest repository work includes assembled source-to-ratio-to-finalizer
bridges, a generic-ring evaluator, 31-to-16 elimination, and arithmetic
source-equivalence work for an actual 89-record source. This progress
must not be conflated with a certificate for a new literal circuit.
The new packet's adapted formula identities and global endpoint still
require the corresponding elaboration and proof-kernel replay. The
pinned formal status documents are retained as inert evidence.

\paragraph{Why the fixed integers are not a loophole.}
The coefficient tuple is produced from a finite, explicitly specified
proof-search source, not selected according to its unknown halting
behavior. Payload hashes bind exact integers to exact source conventions.
A large coefficient is counted as a fixed leaf in the operation ledger,
not pretended to be free of specification cost. Raw and compressed sizes
are separately reported. Different compiler-port conventions, especially
shifted and native masks, are never silently interchanged.

\paragraph{Next decisive experiment.}
The most promising departure from the present finalizer is to elaborate
\texttt{SharedAffineProductEquation.lean}, prove its local source bridge,
then materialize and count the global polynomial produced by the direct
bounded-forall compiler. That would test whether replacing ordered
execution by commutative proof balance helps after compilation, rather
than only at the local level. The common multiple, sequence CRT scale,
and internal compiler majorants must remain separate until an explicit
noncircular sharing theorem is proved. A local 17- or 23-witness count
alone says nothing about the eventual nine- or ten-variable frontier.

A second experiment is integrating the asymmetric parity theorem with
a low-cost binary-base guard. The missing binary-base implication is a
specific obstacle, not a generic appeal to DPRM. Neither experiment is
reported as a completed smaller equation here.

\section{References and retained evidence}
\begin{enumerate}
\item ProveIt, main commit
\texttt{9cca523d9f27ad39d52993f4065cbd7e58f8bdd4},
\url{https://github.com/VladimirReshetnikov/ProveIt}.
The eight reviewed package identities, archive Git blobs and retained
members are recorded in the accompanying source ledgers.
\item Z.-W. Sun, \emph{Further results on Hilbert's Tenth Problem},
Theorems4.12--4.13,
\url{https://arxiv.org/html/1704.03504v7}.
\item Z.-W. Sun, \emph{Ten unknowns for Hilbert's tenth problem over the
integers}, modified Lucas certificate and interpolation,
Theorems3.1 and4.1 and Propositions6.1--6.2,
\url{https://arxiv.org/html/2609.01594v1}.
The exact hypotheses used by the new slice argument are stated in
Section\,\ref{sec:free-lucas-coordinate}; the audit provides a separate proof file.
\item J. P. Jones, D. Sato, H. Wada and D. Wiens,
\emph{Diophantine representation of the set of prime numbers},
American Mathematical Monthly83 (1976),449--464, especially Definition3.7,
Lemma3.8 and Section3.9.
\url{https://www.cs.umd.edu/~gasarch/BLOGPAPERS/Jonesh10.pdf}.
The complete copyrighted article is not redistributed in this archive.
\item The inherited C222 numerical-calculus endpoint and constructive
Diophantine compiler API are retained under
\texttt{direct\_audit/inherited/}, with their own earlier exact repository
pin. These source snapshots expose the interfaces; their transitive
imports are repository dependencies, not newly proved content.
\item Full auxiliary mathematical proofs appear in the component
\texttt{PROOF.md}, \texttt{REVIEW.md}, \texttt{AUDIT.md}, and named theorem
files. Read them together with the coefficient and source manifests.
\end{enumerate}
\paragraph{Reuse.}
Preserved upstream notices identify inherited material. Newly authored
research text, checkers and polynomial data are supplied for incorporation
into ProveIt. The packet asserts neither ownership of inherited material
nor permission to redistribute third-party articles beyond their licenses.

\end{document}
